
\begin{center}\rule{0.5\linewidth}{0.5pt}\end{center}

\section{Статус готовности и сравнение с ТЗ}\label{KT-3.2_Final_Report-ux441ux442ux430ux442ux443ux441-ux433ux43eux442ux43eux432ux43dux43eux441ux442ux438-ux438-ux441ux440ux430ux432ux43dux435ux43dux438ux435-ux441-ux442ux437}

\textbf{MVP-вывод:} проект достиг стадии MVP. Уровень соответствия изначальному ТЗ --- \textbf{\textasciitilde85\%} (upfront от плана).

\textbf{Реализовано (основное):}

\begin{itemize}
\tightlist
\item
  Однопроцессный десктоп-движок реального времени (Windows + Linux), Vulkan 1.4.
\item
  Воксельный мир как источник истины (8×8×8 чанки, dirty-чанки, raycast).
\item
  Greedy-мешинг на GPU (compute shader).
\item
  Walk-персонаж (Jolt \texttt{CharacterVirtual}) с разделением полномочий на запись collision-данных.
\item
  Тени, привязанные к сцене (CSM, оптимизирован под плотные сцены).
\item
  Ray march проход для объёмного рендера.
\item
  TAA (Temporal Anti-Aliasing).
\item
  Клеточный автомат жидкости.
\item
  ECS (flecs) как зеркало для запросов.
\item
  Конвейер ассетов (glTF / Draco / meshopt, горячая перезагрузка).
\item
  5 материалов (Air, Glass, Fluid, FloorWhite, FloorGray) --- см. \texttt{VoxelMaterial} в \texttt{src/voxel/VoxelWorld.hpp}.
\item
  5 пресетов сцен (VoxelLab, EmptyWorld, ShadowTest, FluidDemo).
\item
  Debug HUD, runtime smoke captures, Tracy-профилирование.
\item
  12 ctest suites, \textasciitilde140 тестов, 100\% passing (см. \texttt{docs/KT-2.2\_Test\_Report.md} §Сводный дашборд).
\item
  Аудио (miniaudio) с автообновлением плейлиста.
\item
  Горячая перезагрузка шейдеров через F5.
\item
  Цикл оптимизаций производительности (профилирование + рефакторинг горячих путей).
\end{itemize}

\textbf{Отложено в технический долг:}

\begin{itemize}
\tightlist
\item
  SVO (Sparse Voxel Octree) --- Phase 4 Vision, требует переписывания рендерера под рейтрейсинг (SVO и greedy-мешинг --- разные парадигмы: SVO для ray-marching, greedy для полигонального мешинга).
\item
  Mesh shaders --- не поддерживаются целевой GPU (Vulkan 1.4 fallback на compute).
\item
  Bindless descriptors --- избыточно для текущего размера сцены; планируется в Phase 6 Vision.
\item
  Save/load сериализация мира --- отсутствует (мир создаётся в runtime из пресета).
\end{itemize}

\section{Архитектурные и инженерные вызовы (Post-mortem)}\label{KT-3.2_Final_Report-ux430ux440ux445ux438ux442ux435ux43aux442ux443ux440ux43dux44bux435-ux438-ux438ux43dux436ux435ux43dux435ux440ux43dux44bux435-ux432ux44bux437ux43eux432ux44b-post-mortem}

\paragraph{Post-mortem 1: Walk Authority --- три недели на правильную синхронизацию}\label{KT-3.2_Final_Report-post-mortem-1-walk-authority--ux442ux440ux438-ux43dux435ux434ux435ux43bux438-ux43dux430-ux43fux440ux430ux432ux438ux43bux44cux43dux443ux44e-ux441ux438ux43dux445ux440ux43eux43dux438ux437ux430ux446ux438ux44e}

\textbf{Симптомы:} на ранних этапах walk-персонаж периодически «проскакивал» через стены или залипал в геометрии. Воспроизводимость --- 5--10\% сессий, в зависимости от сложности сцены.

\textbf{Расследование:} Jolt \texttt{CharacterVirtual} имел собственный collision query, независимый от \texttt{VoxelWorld}. Две системы могли расходиться: пока \texttt{VoxelWorld} обновлялся, Jolt продолжал использовать устаревший snapshot. На быстром движении (\textgreater{} 30 m/s) персонаж в одном кадре оказывался «в стене» с точки зрения Jolt, а в следующем кадре \texttt{VoxelWorld} уже обновлялся, и collision query давал другой результат.

\textbf{Решение:} введена концепция \texttt{editVersion} в \texttt{VoxelWorld} (инкремент при любом изменении). \texttt{CharacterVirtual} синхронизируется с \texttt{editVersion} перед каждым physics step. Полномочия на collision query --- только у физики (\texttt{PhysicsWorld}), \texttt{VoxelWorld} --- read-only для других подсистем. Это решило проблему на 100\%, но заняло 3 недели (планировалось 3 дня).

\textbf{Урок:} любая интеграция нативной C++-библиотеки (Jolt, draco) с собственной моделью consistency требует явной фиксации «кто имеет полномочия» (ownership). Без этого --- race conditions в самых неожиданных местах.

\paragraph{Post-mortem 2: SIMD-рефакторинг типов и падение тестов}\label{KT-3.2_Final_Report-post-mortem-2-simd-ux440ux435ux444ux430ux43aux442ux43eux440ux438ux43dux433-ux442ux438ux43fux43eux432-ux438-ux43fux430ux434ux435ux43dux438ux435-ux442ux435ux441ux442ux43eux432}

\textbf{Симптомы:} после серии изменений в \texttt{core/Types.hpp} (Vec3 → SIMD-вектор) все 157 тестов в \texttt{VoxelWorldTests} начали падать с segmentation fault.

\textbf{Расследование:} изменился внутренний layout типа \texttt{Vec3}, но несколько файлов, использующих его через \texttt{inline}-функции в других заголовках, продолжали компилироваться со старыми смещениями полей. ABI несовместимо.

\textbf{Решение:} пересобраны все \texttt{inline}-функции, зависящие от \texttt{Vec3}, в единый header-only модуль; добавлены compile-time проверки через \texttt{static\_assert} (размер \texttt{sizeof(Vec3)}) в каждом потребителе. Все 157 тестов восстановлены.

\textbf{Урок:} при изменении layout типа, на который ссылаются \texttt{inline}-функции в других translation units, нужна полная пересборка всех зависимых единиц. Compile-time проверки (\texttt{static\_assert}) в публичных заголовках --- дешёвый способ ловить такие несовместимости на этапе компиляции.

\paragraph{Post-mortem 3: VoxelLab tremor + TAA descriptor race}\label{KT-3.2_Final_Report-post-mortem-3-voxellab-tremor--taa-descriptor-race}

\textbf{Симптомы:} на статичной сцене VoxelLab (без движения камеры) меши слегка дрожат при включённом TAA. Sidecar-метаданные показывают jittering в \texttt{frame\_time} (±0.3 ms).

\textbf{Расследование:} \texttt{vkAllocateDescriptorSets} для TAA pass вызывался без синхронизации с \texttt{vkQueueSubmit} предыдущего кадра. Descriptor pool исчерпывался под нагрузкой.

\textbf{Решение:} \texttt{vkWaitForFences} с 10 ms timeout перед \texttt{vkAllocateDescriptorSets}; перевод на per-frame SSBO double-buffer (\texttt{SceneFrameResources{[}2{]}}). Частично помогло --- tremor уменьшился, но не пропал полностью. Bug остаётся открытым (BUG-004).

\textbf{Урок:} Vulkan descriptor lifetime --- сложная тема; интегрировать с TAA (которая и так требует frame-coherent state) --- особенно рискованно. Требуется фундаментальный рефакторинг TAA pass.

\section{Организация рабочего процесса}\label{KT-3.2_Final_Report-ux43eux440ux433ux430ux43dux438ux437ux430ux446ux438ux44f-ux440ux430ux431ux43eux447ux435ux433ux43e-ux43fux440ux43eux446ux435ux441ux441ux430}

\textbf{Команда:} 6 человек. Координация --- через \texttt{TODO.md} (чеклисты + roadmap) + \texttt{docs/} (архитектурные ADR и post-mortem) + еженедельные sync-сессии (1 час, аудио + screen-share).

\textbf{Git workflow:}

\begin{itemize}
\tightlist
\item
  \textbf{Основная ветка:} \texttt{master} (всегда green, ctest 12/12 passing).
\item
  \textbf{Feature branches:} \texttt{feature/\textless{}name\textgreater{}} для подзадач; интеграция через pull-request review (минимум 1 approve от участника, не являющегося автором).
\item
  \textbf{Теги:} версионные (например, \texttt{v0.1-mvp}).
\end{itemize}

\textbf{Таск-трекинг:} чекбоксы в \texttt{TODO.md} + краткий daily-standup (15 мин, текстом в общем чате). Не используется Trello/Jira (overhead для команды из 6 человек --- проще держать roadmap в репозитории рядом с кодом).

\textbf{Цикл оптимизаций (закрытый):} серия профилировочных итераций --- от простых (branch hints) до сложных (kernel рефакторинг). Каждая итерация: замер Tracy → профилирование → правка → замер. Все итерации --- закрытые, baseline производительности вырос на каждом шаге.

\section{Рефлексия и приобретённый опыт}\label{KT-3.2_Final_Report-ux440ux435ux444ux43bux435ux43aux441ux438ux44f-ux438-ux43fux440ux438ux43eux431ux440ux435ux442ux451ux43dux43dux44bux439-ux43eux43fux44bux442}

\textbf{7 ключевых уроков:}

\begin{enumerate}
\def\labelenumi{\arabic{enumi}.}
\tightlist
\item
  \textbf{Walk authority (см. Post-mortem 1)} --- любая интеграция нативной C++-библиотеки требует явной фиксации ownership данных.
\item
  \textbf{Layout типы + inline (см. Post-mortem 2)} --- \texttt{static\_assert(sizeof(...))} в публичных заголовках ловит ABI-несовместимости на этапе компиляции.
\item
  \textbf{TAA descriptor race (см. Post-mortem 3)} --- Vulkan descriptor pool исчерпывается под нагрузкой; нужен explicit fence synchronization.
\item
  \textbf{Vulkan 1.4 ≠ магия} --- каждая новая фича (dynamic rendering, timeline semaphores) требует ручной валидации; полагаться на validation layers как «source of truth» --- ошибка.
\item
  \textbf{Greedy vs SVO} --- на маленьких сценах (\textless{} 10K вокселей) greedy-мешинг выигрывает по простоте; SVO оправдан только на больших сценах.
\item
  \textbf{Smoke-тесты должны быть автоматическими} --- ручной запуск \texttt{Invoke-ProjectVRuntimeSmoke.sh} (Linux) / \texttt{Invoke-ProjectVRuntimeSmoke.ps1} (Windows) тратит время; нужен CI workflow.
\item
  \textbf{Инкрементальные коммиты экономят время} --- большие «коммиты-пакеты» трудно откатывать и bisect\textquotesingle ить.
\end{enumerate}

\textbf{Что сделал бы иначе:}

\begin{itemize}
\tightlist
\item
  Сразу ввести \texttt{editVersion} в \texttt{VoxelWorld} (а не после того, как walk authority начал ломаться).
\item
  Использовать SVO вместо greedy-мешинга с самого начала (SVO лучше масштабируется).
\item
  Автоматизировать CI smoke-тесты для runtime captures (сейчас ручной запуск \texttt{Invoke-ProjectVRuntimeSmoke.sh} / \texttt{Invoke-ProjectVRuntimeSmoke.ps1}).
\item
  Делать инкрементальные \texttt{git\ commit} каждый день (а не накапливать изменения).
\end{itemize}

\textbf{Новые изученные паттерны/фреймворки/алгоритмы:}

\begin{itemize}
\tightlist
\item
  Flecs ECS (RAII, MIT, header-only).
\item
  Jolt Physics (\texttt{CharacterVirtual}, deterministic).
\item
  Vulkan 1.4 dynamic rendering + timeline semaphores.
\item
  Greedy-мешинг (Mikola Lysenko, 2012, блог "0 FPS": \url{https://blog.0fps.net/2013/09/25/ambient-occlusion-for-minecraft-like-worlds/}).
\item
  Sparse Voxel Octree (Laine \& Karras, 2010) --- не реализован, но изучен.
\end{itemize}

\section{План дальнейшего развития (Roadmap)}\label{KT-3.2_Final_Report-ux43fux43bux430ux43d-ux434ux430ux43bux44cux43dux435ux439ux448ux435ux433ux43e-ux440ux430ux437ux432ux438ux442ux438ux44f-roadmap}

\begin{longtable}[]{@{}llll@{}}
\toprule\noalign{}
Фаза & Срок & Что & Цель \\
\midrule\noalign{}
\endhead
\bottomrule\noalign{}
\endlastfoot
Phase 4 (Vision) & Q3 2026 & SVO (Sparse Voxel Octree) & Масштабирование на 1M+ вокселей \\
Phase 5 & Q3 2026 & Bindless descriptors & Упрощение renderer, меньше binding-switching \\
Phase 6 & Q4 2026 & Mesh shaders (где поддерживается) & Ускорение greedy-мешинга в 2-3× \\
Phase 7 & Q4 2026 & Save/load сериализация мира & Persistence между запусками \\
Phase 8 & Q1 2027 & Сетевая игра & Collaborative editing воксельного мира \\
Phase 9 & Q2 2027 & WebAssembly build (WASM) & Запуск в браузере (Emscripten + Vulkan-через-WebGPU) \\
\end{longtable}

